% Compile with LuaLaTeX so that the Unicode in the Lean listings renders.
\documentclass{article}
\usepackage{amsmath, amssymb, amsthm}
\usepackage{hyperref}
\usepackage[margin=1in]{geometry}
\usepackage{fontspec}
\usepackage{xcolor}
\usepackage{listings}
\usepackage{needspace}
\usepackage{titlesec}

% Unicode-capable monospace font for the Lean code (has ℂ ℝ ℚ ∀ ∃ ⟨ ⟩ ∥ ⟂ ⊔ ⊥ ↔ ↥ ...)
\IfFontExistsTF{DejaVu Sans Mono}{\setmonofont{DejaVu Sans Mono}[Scale=0.88]}
{\IfFontExistsTF{Menlo}{\setmonofont{Menlo}[Scale=0.88]}{\setmonofont{FreeMono}[Scale=0.95]}}
% DejaVu Sans Mono lacks U+2225 (∥): typeset it from a font that has it
\newfontfamily\fbsym{FreeMono}[Scale=0.95]
% Under LuaLaTeX, listings typesets unknown non-ASCII characters outside its token buffer, which
% reorders text next to them (e.g. x⁻¹ or {θ : ℝ}). Route every non-ASCII character used in the
% code through listings' own machinery instead. (∥ is missing from DejaVu Sans Mono: use FreeMono.)
\makeatletter
\catcode`α=\active \catcode`θ=\active \catcode`⁻=\active \catcode`ℂ=\active \catcode`ℕ=\active \catcode`ℚ=\active \catcode`ℝ=\active \catcode`→=\active \catcode`↔=\active \catcode`↥=\active \catcode`∀=\active \catcode`∃=\active \catcode`∈=\active \catcode`∉=\active \catcode`∣=\active \catcode`∥=\active \catcode`∧=\active \catcode`∨=\active \catcode`≠=\active \catcode`⊔=\active \catcode`⊥=\active \catcode`⟂=\active \catcode`⟨=\active \catcode`⟩=\active
\defα{\lst@ProcessOther{{\ttfamily\char945 }}}
\defθ{\lst@ProcessOther{{\ttfamily\char952 }}}
\def⁻{\lst@ProcessOther{{\ttfamily\char8315 }}}
\defℂ{\lst@ProcessOther{{\ttfamily\char8450 }}}
\defℕ{\lst@ProcessOther{{\ttfamily\char8469 }}}
\defℚ{\lst@ProcessOther{{\ttfamily\char8474 }}}
\defℝ{\lst@ProcessOther{{\ttfamily\char8477 }}}
\def→{\lst@ProcessOther{{\ttfamily\char8594 }}}
\def↔{\lst@ProcessOther{{\ttfamily\char8596 }}}
\def↥{\lst@ProcessOther{{\ttfamily\char8613 }}}
\def∀{\lst@ProcessOther{{\ttfamily\char8704 }}}
\def∃{\lst@ProcessOther{{\ttfamily\char8707 }}}
\def∈{\lst@ProcessOther{{\ttfamily\char8712 }}}
\def∉{\lst@ProcessOther{{\ttfamily\char8713 }}}
\def∣{\lst@ProcessOther{{\ttfamily\char8739 }}}
\def∥{\lst@ProcessOther{{\fbsym\char8741 }}}
\def∧{\lst@ProcessOther{{\ttfamily\char8743 }}}
\def∨{\lst@ProcessOther{{\ttfamily\char8744 }}}
\def≠{\lst@ProcessOther{{\ttfamily\char8800 }}}
\def⊔{\lst@ProcessOther{{\ensuremath{\sqcup}}}}
\def⊥{\lst@ProcessOther{{\ensuremath{\bot}}}}
\def⟂{\lst@ProcessOther{{\ensuremath{\perp}}}}
\def⟨{\lst@ProcessOther{{\ttfamily\char10216 }}}
\def⟩{\lst@ProcessOther{{\ttfamily\char10217 }}}
\makeatother



\newtheorem{theorem}{Theorem}

% ---------- inline Lean names (underscores allowed) ----------
\newcommand{\lean}[1]{\texttt{\detokenize{#1}}}

% ---------- Lean listings ----------
\definecolor{leankw}{HTML}{1F4FBF}     % declarations / keywords
\definecolor{leantac}{HTML}{8E24AA}    % tactics
\definecolor{leantype}{HTML}{C2570C}   % core types / Prop
\definecolor{leancom}{HTML}{2E7D32}    % comments
\definecolor{leanstr}{HTML}{B71C1C}    % strings
\definecolor{leanbg}{HTML}{F6F7FB}     % background
\definecolor{leanrule}{HTML}{9FA8DA}   % frame

\lstdefinelanguage{Lean}{
    morekeywords=[1]{theorem,lemma,def,inductive,structure,mutual,end,namespace,open,import,
        instance,where,noncomputable,macro,macro_rules,syntax,elab,elab_rules,section,variable,
        example,by,fun,have,show,from,at,with,in,let,if,then,else,match,using,set,this},
    morekeywords=[2]{intro,intros,exact,apply,rw,simp,constructor,refine,use,obtain,rcases,
        cases,induction,linarith,nlinarith,ring,ring_nf,field_simp,grind,positivity,
        linear_combination,by_cases,by_contra,unfold,exact_mod_cast,omega,first,try,repeat,
        assumption,left,right,rfl,norm_num,subst,simpa,tactic},
    morekeywords=[3]{Prop,Type,Sort,Complex,Real,Line},
    sensitive=true,
    morecomment=[l]{--},
    morecomment=[s]{/-}{-/},
    morestring=[b]",
}

\lstset{
    language=Lean,
    basicstyle=\ttfamily\footnotesize,
    keywordstyle=[1]\color{leankw}\bfseries,
    keywordstyle=[2]\color{leantac},
    keywordstyle=[3]\color{leantype},
    commentstyle=\color{leancom}\itshape,
    stringstyle=\color{leanstr},
    backgroundcolor=\color{leanbg},
    frame=leftline,
    framerule=2pt,
    rulecolor=\color{leanrule},
    xleftmargin=6pt,
    framexleftmargin=4pt,
    columns=fixed,
    keepspaces=true,
    showstringspaces=false,
    breaklines=true,
    breakatwhitespace=true,
    tabsize=2,
    aboveskip=0.8em,
    belowskip=0.6em,
    captionpos=b,
}

\title{\vspace{-2em}Formalizing Origami Geometry in Lean 4}
\author{Manya Kansal}
\date{}

\begin{document}
\maketitle
\vspace{-1em}
% \tableofcontents
\section{Introduction}

The Huzita-Justin axioms (HJAs) formalise the seven fundamental origami fold operations.
Unlike just compass-and-straightedge operations, origami operations can solve cubic
equations via axiom O6 (the Beloch fold) \cite{alperin2000}, enabling constructions such as
angle trisection. The field of origami-constructible numbers is characterised by Alperin's
theorem: $z \in \mathbb{R}$ is origami-constructible iff it lies in an extension of
$\mathbb{Q}$ of degree $2^a3^b$. Despite the elegance of these results,
no machine-verified proof exists (that I could find).

This report gives an overview of the problem tackled (formalising origami constructibility in
Lean 4 with Mathlib), the theorems proved, the techniques used, and what I learnt about Lean
that was not covered in class.

\section{Original Goals}

\begin{enumerate}
    \item Define terms for folds (\texttt{FoldLine}, \texttt{foldPoint},
    \texttt{foldLine}) and the HJAs O1--O7 as a synthetic axiomatic
    system rather than a coordinate system, similar to the SystemE/LeanGeo library \cite{avigad2009}.

    \item Prove basic fold properties: isometry, injectivity, and that
    the O2 fold line is the perpendicular bisector of the two given points, etc.

    \item Formalize \textbf{Abe's angle trisection} \cite{fuchs2011}:
    express the construction in terms of HJA axioms, then prove
    correctness by reducing to the identity
    $\cos(3\varphi) = 4\cos^3(\varphi) - 3\cos(\varphi)$
    using Mathlib's trigonometry library.

    \item If time permits, characterize which regular $n$-gons are
    origami-constructible (those with $\phi(n) = 2^c3^d$).
\end{enumerate}

\section{Achieved Goals}

\begin{enumerate}
    \item \textbf{Done.} Lines are a point plus a nonzero direction in $\mathbb{C}$ (\lean{Line});
    a fold is the reflection \lean{Line.reflect}; O1--O7 are predicates; constructible points and
    lines form a mutual inductive type. This is an inductive system inside $\mathbb{C}$, not a
    SystemE-style synthetic one.

    \item \textbf{Done.} \lean{reflect_dist_sq} (isometry), \lean{reflect_reflect} (involution,
    hence injective), \lean{reflect_eq_self_iff} (fixed points are exactly the crease),
    \lean{reflect_bisector} and \lean{o2_eq} (the O2 fold line is the perpendicular bisector).

    \item \textbf{Done.} Proved the Beloch fold (\lean{beloch_fold}), that O6 constructs
    every real root of a cubic with constructible coefficients (\lean{cpt_cubic}), and the
    trisection identity (\lean{trisect_cos}). Abe's fold sequence itself is not written out as an
    explicit HJA construction.

    \item \textbf{Partly done.} The characterisation of regular $n$-gons itself is not proved,
    but the main field-theoretic result it rests on is: \lean{OCField} is a field,
    \lean{alperin} characterises the constructible numbers as those lying in a cubic tower, and
    \lean{origami_degree} shows that their degree is of the form $2^a3^b$.

    \item Beyond the original goals, I proved the algebra--geometry equivalence (soundness and completeness).
\end{enumerate}

\section{Design Decisions}
\label{sec:design}

\subsection{Why Complex Coordinates}
I initially defined all the predicates synthetically, similar to \texttt{LeanGeo}.
However, that would need too many axioms to be stated by hand, possibly leading to
contradictions. Also, it was harder to write general tactics for synthetic proofs. As a result,
complex coordinates are used.

\subsection{Membership Through Conjugation}
A point $z$ lies on the line through $p$ with direction $d$ exactly when $z-p$ is a real multiple
of $d$, i.e.\ when $(z-p)\,\overline{d}$ is a real number. This is why membership is defined as
\lean{((z - l.p) * conj l.d).im = 0}, which is equivalent to:
\begin{lstlisting}
theorem mem_iff_conj (l : Line) (z : ℂ) :
    z ∈ l ↔ (z - l.p) * conj l.d = conj (z - l.p) * l.d := by ...
\end{lstlisting}
It needs no division and no special case for vertical lines, unlike writing a line as 
$y=mx+c$ or solving for a parameter $t$ in $z=p+td$. Parallelism and perpendicularity use 
the same idea: \lean{Par} asks that $d\,\overline{d'}$ is real
and \lean{Perp} asks that it is purely imaginary.

\subsection{Why a Line Carries a Proof That Its Direction is Nonzero}
If $d=0$, then $(z-p)\,\overline{d}=0$ for every $z$, so the ``line'' would contain every point of
the plane. Also, \lean{reflect} divides by \lean{conj d}. Putting the proof \lean{hd} inside 
the structure means that no theorem needs \lean{d ≠ 0} as an extra hypothesis, and lemmas 
such as \lean{Line.d_re_or_im} can use it directly.

\subsection{Why \texttt{Same} is a Setoid}
The same geometric line has many representations: any point on it can be the base point $p$, and
$d$ can be scaled by a nonzero real number. \lean{Same l m} says that $l$ and $m$ are parallel and
that \lean{l.p} lies on $m$. This relation is reflexive, symmetric and transitive
(\lean{par_refl} and \lean{mem_self}, \lean{same_symm}, \lean{same_trans}), so it is an
equivalence relation, and it is registered as a \lean{Setoid} instance. It is also the reason
\lean{GeoConstrLine} has the constructor \lean{same}: constructibility must not depend on which
representation of a line is used. The following proofs benefit from this:
\begin{lstlisting}
theorem par_refl (l : Line) : geo{l ∥ l} := by ...
theorem same_symm {l m : Line} (h : l.Same m) : m.Same l := by ...
theorem same_trans {l m n : Line} (h1 : l.Same m) (h2 : m.Same n) : l.Same n := by ...
\end{lstlisting}

\section{The Structure of the Project}
\label{sec:structure}
Points are complex numbers and lines are defined by a point (a complex number which it passes
through) and a direction (a nonzero complex number in the direction of the line), along with a
proof that the direction is nonzero. Membership, parallelism and reflection are all expressed
with complex conjugation.
\begin{lstlisting}
structure Line where
    p : ℂ
    d : ℂ
    hd : d ≠ 0

instance : Membership ℂ Line where
  mem l z := ((z - l.p) * conj l.d).im = 0

def Par  (l m : Line) : Prop := (l.d * conj m.d).im = 0
def Perp (l m : Line) : Prop := (l.d * conj m.d).re = 0
def Same (l m : Line) : Prop := Par l m ∧ l.p ∈ m
noncomputable def reflect (f : Line) (z : ℂ) : ℂ :=
  f.p + (f.d / conj f.d) * conj (z - f.p)
\end{lstlisting}
There are 2 components to the project: the algebra-geometry equivalence
and the algebraic characterisation of origami-constructible numbers.

\subsection{The Algebra-Geometry Equivalence}
\label{sec:alg-geo-eq}
A point is Geometrically constructible if it is \(0\), \(1\), an intersection of two origami-constructible
lines, or the reflection of an origami-constructible point in an origami-constructible line.
A line is origami-constructible if it is constructible by one of the \(7\) HJAs.
The \texttt{geo\{\dots\}} notation is a small custom syntax category for $\parallel$, $\perp$
and $\equiv$. In Lean, I defined these as:
\begin{lstlisting}
def O1 (p q : ℂ) (f : Line) : Prop := p ∈ f ∧ q ∈ f
def O2 (p q : ℂ) (f : Line) : Prop := f.reflect p = q
def O3 (l m f : Line) : Prop := (f.reflectLine l).Same m
def O4 (p : ℂ) (l f : Line) : Prop := geo{ f ⟂ l ∧ p ∈ f }
def O5 (p q : ℂ) (l f : Line) : Prop := f.reflect p ∈ l ∧ q ∈ f
def O6 (p q : ℂ) (l m f : Line) : Prop := f.reflect p ∈ l ∧ f.reflect q ∈ m
def O7 (p : ℂ) (l m f : Line) : Prop := geo{ f ⟂ m ∧ (f.reflect p) ∈ l }

mutual
inductive GeoConstrPt : ℂ → Prop
| zero : GeoConstrPt 0
| one  : GeoConstrPt 1
| inter {l m : Line} {z : ℂ} : GeoConstrLine l → GeoConstrLine m → ¬ l.Par m → z ∈ l → z ∈ m → GeoConstrPt z
| reflect {p : ℂ} {f : Line} : GeoConstrPt p → GeoConstrLine f → GeoConstrPt (f.reflect p)

inductive GeoConstrLine : Line → Prop
| o1 {p q : ℂ} {f : Line} :         -- line through p and q
    GeoConstrPt p → GeoConstrPt q → p ≠ q → O1 p q f → GeoConstrLine f
| o2 {p q : ℂ} {f : Line} :         -- line sending p to q
    GeoConstrPt p → GeoConstrPt q → p ≠ q → O2 p q f → GeoConstrLine f
| o3 {l m f : Line} :               -- angle bisector
    GeoConstrLine l → GeoConstrLine m → ¬ l.Same m → O3 l m f → GeoConstrLine f
| o4 {p : ℂ} {l f : Line} :         -- line perpendicular to l through p
    GeoConstrPt p → GeoConstrLine l → O4 p l f → GeoConstrLine f
| o5 {p q : ℂ} {l f : Line} :       -- line sending p to l, passing through q
    GeoConstrPt p → GeoConstrPt q → GeoConstrLine l → p ≠ q → O5 p q l f → GeoConstrLine f
| o6 {p q : ℂ} {l m f : Line} :     -- line sending p to l and q to m
    GeoConstrPt p → GeoConstrPt q → GeoConstrLine l → GeoConstrLine m → p ∉ l → ¬ geo{l ∥ m} →
    O6 p q l m f → GeoConstrLine f
| o7 {p : ℂ} {l m f : Line} :       -- line sending p to l, perpendicular to m
    GeoConstrPt p → GeoConstrLine l → GeoConstrLine m → ¬ geo{l ∥ m} → O7 p l m f → GeoConstrLine f
| same {f g : Line} : GeoConstrLine f → f.Same g → GeoConstrLine g
end
\end{lstlisting}
A number is Algebraically constructible if it lies in the set of numbers containing \(0\) and \(1\),
which is closed under addition, multiplication, additive and multiplicative inverses, and taking
square roots and cube roots. In Lean, I defined it as:
\begin{lstlisting}
inductive AlgConstrNum : ℝ → Prop
| rat (q : ℚ) : AlgConstrNum q
| add {x y} : AlgConstrNum x → AlgConstrNum y → AlgConstrNum (x + y)
| mul {x y} : AlgConstrNum x → AlgConstrNum y → AlgConstrNum (x * y)
| neg {x} : AlgConstrNum x → AlgConstrNum (-x)
| inv {x} : AlgConstrNum x → AlgConstrNum x⁻¹
| sqrt {x} : AlgConstrNum x → AlgConstrNum (Real.sqrt x)
| cubic {a b c r : ℝ} : AlgConstrNum a → AlgConstrNum b → AlgConstrNum c →
    r ^ 3 + a * r ^ 2 + b * r + c = 0 → AlgConstrNum r

def AlgConstrLine (f : Line) : Prop :=
    ∃ a b c : ℝ, AlgConstrNum a ∧ AlgConstrNum b ∧ AlgConstrNum c ∧ (a ≠ 0 ∨ b ≠ 0) ∧
        ∀ z : ℂ, z ∈ f ↔ a * z.re + b * z.im + c = 0
\end{lstlisting}
The corresponding equivalence theorems are:
\begin{lstlisting}
theorem GeoConstrPt_iff (z : ℂ) : GeoConstrPt z ↔ AlgConstrNum z.re ∧ AlgConstrNum z.im :=
    ⟨sound_pt, complete_pt⟩

theorem GeoConstrLine_iff (f : Line) : GeoConstrLine f ↔ AlgConstrLine f :=
    ⟨sound_line, complete_line⟩
\end{lstlisting}

\subsubsection{Soundness (geometric $\Rightarrow$ algebraic)}
\label{sec:soundness}
Each axiom is translated into coordinates, giving one lemma per axiom (\lean{o1_aux}, \dots,
\lean{o7_aux}). The proof uses \texttt{mutual} recursion, similar to its \texttt{mutual} definition.
\begin{lstlisting}
mutual
theorem sound_pt : ∀ {z : ℂ}, GeoConstrPt z → AlgConstrNum z.re ∧ AlgConstrNum z.im
| _, .zero => by simp [AlgConstrNum.zero]
| _, .one => by simp [AlgConstrNum.one, AlgConstrNum.zero]
| _, .inter hl hm hn hzl hzm => inter_aux (sound_line hl) (sound_line hm) hn hzl hzm
| _, .reflect hp hf => reflect_aux (sound_pt hp) (sound_line hf)

theorem sound_line : ∀ {f : Line}, GeoConstrLine f → AlgConstrLine f
| _, .o1 hp hq hne h => o1_aux (sound_pt hp) (sound_pt hq) hne h
-- o2 ... o6 are analogous
| _, .o7 hp hl hm hlm h => o7_aux (sound_pt hp) (sound_line hl) (sound_line hm) hlm h
| _, .same hf he => by ...
end
\end{lstlisting}

\subsubsection{Completeness (algebraic $\Rightarrow$ geometric)}
I first construct the axes and the operations $+,-,\times,{}^{-1}$ (multiplicative inverse) on the $x$-axis by folds
(\lean{cpt_add}, \lean{cpt_mul}, \lean{cpt_inv}), then \lean{cpt_mk} builds $u+vi$ from $u$
and $v$. Cubic roots use the Beloch fold, and square roots come for free because $\sqrt a$ is a
root of $t^3-at=0$.
\begin{lstlisting}
theorem beloch_fold {t a b c : ℝ} (h : t ^ 3 + a * t ^ 2 + b * t + c = 0)
    (T : ℂ) (hTre : T.re = 1 - b) (hTim : T.im = a - c) :
    ((bfold t).reflect 0).re = 2 ∧ ((bfold t).reflect T).im = a + c := by ...

theorem cpt_cubic {a b c r : ℝ} (ha : RC a) (hb : RC b) (hc : RC c)
    (h : r ^ 3 + a * r ^ 2 + b * r + c = 0) : GeoConstrPt r := by ...

theorem cpt_sqrt {a : ℝ} (ha : RC a) : RC (Real.sqrt a) := by ...
\end{lstlisting}
\subsection{The Algebraic Characterisation}
Alperin's Theorem~\cite{alperin2000,lee2017} states that a real number is origami-constructible if
and only if it lies in an algebraic extension of $\mathbb{Q}$ built by adjoining roots of cubics.
The constructible numbers form a subfield \lean{OCField} of $\mathbb{R}$ containing $\mathbb{Q}$.
\lean{Reach K} says $K$ is built from $\mathbb{Q}$ by adjoining roots of cubics one at a time,
and \lean{alperin} shows a number is constructible iff it lies in such a $K$.
Each step of a \lean{Reach} tower has degree 1, 2 or 3 (\lean{finrank_step}, via the tower law), so
$[\mathbb{Q}(x):\mathbb{Q}]$ is of the form $2^a3^b$ (\lean{origami_degree}).
\begin{lstlisting}
inductive Reach : IntermediateField ℚ ℝ → Prop
| bot : Reach ⊥
| step {K : IntermediateField ℚ ℝ} {α a b c : ℝ} :
    Reach K → a ∈ K → b ∈ K → c ∈ K →
    α ^ 3 + a * α ^ 2 + b * α + c = 0 →
    Reach (K ⊔ IntermediateField.adjoin ℚ {α})

theorem alperin {x : ℝ} : AlgConstrNum x ↔ ∃ K, Reach K ∧ x ∈ K := by ...

def isTwoThree (n : ℕ) : Prop := ∀ p, p.Prime → p|n → p = 2 ∨ p = 3

theorem origami_degree {x : ℝ} (h : AlgConstrNum x) :
    isTwoThree (Module.finrank ℚ ↥(IntermediateField.adjoin ℚ {x})) := by ...
\end{lstlisting}
\subsection{Angle Trisection}
If $\cos\theta$ is constructible then so is $\cos(\theta/3)$. Mathlib's \lean{Real.cos_three_mul}
shows $x=\cos(\theta/3)$ is a root of $x^3-\tfrac34x-\tfrac14\cos\theta=0$, so the cubic rule
applies.
\begin{lstlisting}
theorem trisect_cos {θ : ℝ} (h : AlgConstrNum (Real.cos θ)) : AlgConstrNum (Real.cos (θ / 3)) := by
    have h3 := Real.cos_three_mul (θ / 3)
    have : 3 * (θ / 3) = θ := by linarith
    rw [this] at h3
    refine AlgConstrNum.cubic (a := 0) (b := -3 / 4) (c := -Real.cos θ / 4)
        AlgConstrNum.zero (by is_oc) (by is_oc) ?_
    linear_combination (-1 / 4) * h3
\end{lstlisting}

\section{Techniques Used}
\begin{itemize}
    \item \textbf{Reduction to coordinates:} rewriting $\mathbb{C}$-valued statements into real and
    imaginary parts (\lean{to_re_im}) so that real arithmetic closes the goal. Conjugation is
    handled separately by \lean{conj_simp}, which pushes \lean{conj} through sums, products and
    quotients, so that many goals can be closed without leaving $\mathbb{C}$ at all.
    
    \item \textbf{Algebraic closers:} \lean{linear_combination}, \lean{field_simp},
    \lean{positivity}, \lean{nlinarith}, \lean{grind}. Most geometric facts become polynomial
    identities once the coordinates are expanded. \lean{field_simp} clears denominators (using
    the nonzero-direction facts), and \lean{ring}, \lean{linear_combination} or \lean{nlinarith}
    finish the goal; \lean{grind} handles the case splits and the small inequalities.
    
    \item \textbf{Custom tactics:} \lean{is_oc} and \lean{rc} (closure side-goals),
    \lean{line_facts}, \lean{cancel_dir}, \lean{extract_all_points}. The tactic \lean{line_facts}
    is an \lean{elab} that scans the context for every \lean{Line} and adds the facts that are
    needed every time (\lean{d != 0}, \lean{normSq d > 0}, and so on). The tactic
    \lean{line_unfold} unfolds the definitions of membership, \lean{Par}, \lean{Perp},
    \lean{Same} and the reflections. These are combined into one tactic, \lean{line_auto}, which
    tries several pipelines in turn.
    
    \item \textbf{Automation with fallbacks:} \lean{line_auto} and \lean{line_arith} use
    \lean{first | ... | ...}, so one macro covers many different goals, and the first pipeline that
    works is used. This removed most of the repeated boilerplate from the proofs of the
    individual folds.
    
    \item \textbf{Case analysis on directions:} since a \lean{Line} carries a proof that its
    direction is nonzero, either the real or the imaginary part is nonzero. \lean{cancel_dir}
    splits on this and cancels the nonzero factor from both sides of an equation.
    
    \item \textbf{Reduction of O6 to a cubic:} the Beloch fold is handled by showing that a number
    is constructible whenever it is a root of a cubic whose coefficients are constructible
    (\lean{AlgConstrNum.of_cubic_fn}). Both the soundness and the completeness of O6 reduce to
    this one lemma.
    
    \item \textbf{Mathlib field theory:} There are in-built structures and functions \lean{IntermediateField}, 
    \lean{adjoin}, \lean{minpoly} and the tower law \lean{Module.finrank_mul_finrank} for the degree bound. 
    Alperin's necessary condition is proved by induction on the construction: each step adjoins 
    a root of a cubic, so it multiplies the degree by a factor of at most 3 (or exactly 2, for a quadratic).
\end{itemize}

\section{Things that I learnt, that were not taught in class.}

\subsection{The \texttt{refine} Tactic}
\texttt{refine} is similar to \texttt{exact}, but \lean{?_} marks the parts that have not yet
been proved and they become new goals. It lets us supply some arguments of a theorem and turn
the rest into goals. For example, in \lean{o1_aux} the witnesses of the existential goal are
chosen with \lean{use}, and \lean{refine} then proves four of the five conditions on the spot
and leaves only the last one as a new goal:
\begin{lstlisting}
use -(q.im - p.im), q.re - p.re, (q.im - p.im) * p.re - (q.re - p.re) * p.im
refine ⟨by is_oc, by is_oc, by is_oc, by grind, ?_⟩
-- remaining goal: ∀ z : ℂ, z ∈ f ↔ a * z.re + b * z.im + c = 0
\end{lstlisting}
and for lemmas with many arguments, as in \lean{o6_core}, where \lean{refine} names one argument
with \lean{(G := G)}, fills in the easy ones directly, and leaves four goals (the cubic identity
for \lean{G} and the constructibility of \lean{G 0}, \lean{G 1}, \lean{G (-1)}), which are then
closed one by one with `\lean{·}' per goal:
\begin{lstlisting}
have hsI : AlgConstrNum (-b1 * P + a1 * Q) := by
    refine AlgConstrNum.of_cubic_fn (G := G)
        (c3 := (a1 * b2 - a2 * b1) * (a1 ^ 2 + b1 ^ 2) * ((a1 ^ 2 + b1 ^ 2)⁻¹) ^ 3)
        (mul_ne_zero (mul_ne_zero hD hN) (pow_ne_zero 3 hw)) (by is_oc) ?_ ?_ ?_ ?_ hs
    · intro x
        rw [hG x, hG 1, hG (-1), hG 0]
        unfold O6poly
        ring
    · rw [hG]; unfold O6poly; is_oc
    · rw [hG]; unfold O6poly; is_oc
    · rw [hG]; unfold O6poly; is_oc
\end{lstlisting}

\subsection{\texttt{mutual} Recursive Definitions}
Constructible points and lines use each other in their definitions, so they are declared in a
\lean{mutual} block\ref{sec:alg-geo-eq}. Proofs about them\ref{sec:soundness} are also written
as mutual theorems with pattern matching on the constructors; as far as I found,
the \lean{induction} tactic does not handle this case directly.

\subsection{Setoids}
A \lean{Setoid} is a type together with an equivalence relation on it, bundled as one structure:
a relation \lean{r} and a proof \lean{iseqv} that \lean{r} is reflexive, symmetric and transitive.
Once the instance is declared, Lean treats \lean{r} as the default notion of ``being the same'' on
that type: $a \approx b$ means \lean{r a b}, and it becomes possible to form the quotient type,
whose elements are the equivalence classes. In my project, \lean{Line.Same} is the relation,
and the three proofs are those of \lean{par_refl}/\lean{mem_self}, \lean{same_symm} and
\lean{same_trans} (see Section~\ref{sec:design}):
\begin{lstlisting}
instance : Setoid Line where
  r := Line.Same
  iseqv := ⟨fun l => ⟨par_refl l, mem_self l⟩, same_symm, same_trans⟩
\end{lstlisting}

\subsection{Writing My Own Tactics and Syntax}
Tactics allowed cleaner proofs and reduced the amount of hand-computation required:
\lean{macro_rules} can call themselves and recurse through a term to prove \lean{AlgConstrNum}
goals (\lean{is_oc}); \lean{elab} with Lean's metaprogramming loops over the local context
to add facts about every \lean{Line} (\lean{line_facts}). A \lean{declare_syntax_cat} gave
the \lean{geo{...}} notation.
\begin{lstlisting}
syntax "is_oc_core" : tactic
macro_rules
    | `(tactic| is_oc_core) => `(tactic| first
        | assumption
        | exact AlgConstrNum.zero
        | (apply AlgConstrNum.add <;> is_oc_core)
        | (apply AlgConstrNum.mul <;> is_oc_core)
        | (apply AlgConstrNum.sqrt; is_oc_core))

macro "is_oc" : tactic => `(tactic| (try casesm* _ ∧ _) <;> is_oc_core)
\end{lstlisting}

\subsection{Mathlib's Field Theory and Automation}
Working with \lean{IntermediateField} needed \lean{restrictScalars}/\lean{extendScalars} to move
between $\mathbb{Q}$- and $K$-structures. Naming implicit arguments with \lean{@} in an
\lean{induction ... with} case, and building witnesses with
\lean{obtain ⟨x, hx⟩ : ∃ x, ... := ⟨_, _⟩}, were also new to me.

\subsection{How bad Claude is at writing short proofs}
AI suggestions for proofs were often unnecessarily wrong (lemma names that do not exist, tactics
that do not close the goal), so I treated them as hints and tried to write short proofs myself
whenever possible. The funniest example of this was in \lean{o3_eq}, where Claude told me to
first prove the following lemma, relating the sign of $u_1+u_2$ to a sign $s_0$:
\begin{lstlisting}
theorem sign_sum {u1 u2 s0 : ℝ} (hs0 : s0 = 1 ∨ s0 = -1)
    (h1 : (u1 > 0 → s0 = 1) ∧ (u1 < 0 → s0 = -1))
    (h2 : (u2 > 0 → s0 = 1) ∧ (u2 < 0 → s0 = -1))
    (hne : u1 ≠ 0 ∨ u2 ≠ 0) :
    (if u1 + u2 > 0 then (1 : ℝ) else -1) = s0 := by
    rcases hs0 with rfl | rfl
    · have a : 0 ≤ u1 := by
        by_contra h; push_neg at h; have := h1.2 h; norm_num at this
        have b : 0 ≤ u2 := by
        by_contra h; push_neg at h; have := h2.2 h; norm_num at this
        have hsum : 0 < u1 + u2 := by
        rcases hne with h | h
        · have := lt_of_le_of_ne a (Ne.symm h); linarith
        · have := lt_of_le_of_ne b (Ne.symm h); linarith
        rw [if_pos hsum]
    · have a : u1 ≤ 0 := by
        by_contra h; push_neg at h; have := h1.1 h; norm_num at this
        have b : u2 ≤ 0 := by
        by_contra h; push_neg at h; have := h2.1 h; norm_num at this
        have hsum : u1 + u2 < 0 := by
        rcases hne with h | h
        · have := lt_of_le_of_ne a h; linarith
        · have := lt_of_le_of_ne b h; linarith
        rw [if_neg (not_lt.2 hsum.le)]
\end{lstlisting}
whereas the corresponding step in my proof is closed by a single line:
\begin{lstlisting}
· by_cases h : (s1 + s2 > 0) <;> simp [h, s3]
\end{lstlisting}

\begin{thebibliography}{9}
\setlength{\itemsep}{0pt}

\bibitem{avigad2009}
J.~Avigad, E.~Dean, J.~Mumma,
``A formal system for Euclid's Elements,''
\textit{Review of Symbolic Logic}, 2009.

\bibitem{fuchs2011}
C.~Fuchs,
``Angle trisection with origami and related topics,''
\textit{Elemente der Mathematik}, 2011.

\bibitem{alperin2000}
R.~C.~Alperin, R.~J.~Lang,
``One-, two-, and multi-fold origami axioms,''
\textit{Origami}, 2009.

\bibitem{lee2017}
H.~Y.~Lee,
``Origami-constructible numbers,''
Master's thesis, University of Georgia, 2017.

\end{thebibliography}

\end{document}